\documentclass[a4paper,12pt]{article}
\usepackage[T2A]{fontenc}
\usepackage[utf8]{inputenc}
\usepackage[russian]{babel}
\usepackage{amsmath,amssymb,mathtools}
\usepackage{helvet}
\renewcommand{\familydefault}{\sfdefault}
\usepackage{icomma}
\usepackage[margin=22mm,top=23mm,bottom=24mm,headheight=14pt]{geometry}
\usepackage{microtype}
\usepackage{tikz}
\usetikzlibrary{arrows.meta,positioning,shapes.geometric,calc}
\usepackage{tabularx,booktabs,array}
\usepackage{enumitem}
\usepackage{fancyhdr}
\usepackage{hyperref}
\usepackage{setspace}
\usepackage{xcolor}
\usepackage{needspace}
\definecolor{accent}{HTML}{087C80}
\definecolor{darkaccent}{HTML}{145A60}
\definecolor{textgray}{HTML}{42484A}
\definecolor{softgray}{HTML}{6D7778}
\hypersetup{colorlinks=true,linkcolor=darkaccent,urlcolor=darkaccent,pdfauthor={Лёвин Артём Александрович},pdftitle={Чек-лист. Банковские кредиты и проценты. Ксения Клыкова}}
\setstretch{1.18}
\setlength{\parindent}{0pt}
\setlength{\parskip}{6pt}
\setlist[itemize]{leftmargin=*,itemsep=2pt,topsep=4pt}
\setlist[enumerate]{leftmargin=*,itemsep=5pt,topsep=4pt}
\pagestyle{fancy}
\fancyhf{}
\fancyhead[L]{\small\color{textgray}Кредитные задачи и проценты}
\fancyhead[R]{\small\color{textgray}Ксения Клыкова}
\fancyfoot[L]{\footnotesize\color{softgray}Лёвин А. А. \,$\cdot$\, 09.11.2026}
\fancyfoot[R]{\footnotesize\color{softgray}\thepage}
\renewcommand{\headrulewidth}{0pt}
\renewcommand{\footrulewidth}{0pt}
\newcommand{\mainsection}[1]{\vspace{4pt}{\Large\color{darkaccent}\bfseries #1}\par\vspace{3pt}}
\newcommand{\subhead}[1]{\needspace{5\baselineskip}\vspace{7pt}{\large\color{accent}\bfseries #1}\par\vspace{2pt}}
\newcommand{\keyidea}[1]{\textcolor{darkaccent}{\bfseries Ключевая идея.} #1}
\newcommand{\cbox}{\ensuremath{\square}\,}
\newcommand{\ex}[2]{\needspace{8\baselineskip}\par\medskip\noindent{\color{darkaccent}\bfseries #1.}\quad{\bfseries #2}\par\vspace{1pt}}
\newcommand{\ansrow}[2]{#1 & #2\\[2mm]}
\newcommand{\mathnote}[1]{\begin{center}\color{darkaccent}$\displaystyle #1$\end{center}}
\begin{document}
\thispagestyle{empty}
\begin{titlepage}
\centering
\vspace*{13mm}
{\Large\color{darkaccent} МАТЕМАТИКА\par}
\vspace{19mm}
{\LARGE\color{textgray} Чек-лист\par}
\vspace{10mm}
{\LARGE\color{accent} Банковские кредиты:\par}
\vspace{2mm}
{\LARGE\color{accent} проценты и графики погашения\par}
\vspace{12mm}
{\large Подготовка к ЕГЭ по профильной математике\par}
\vfill
{\large Ксения Клыкова\par}
\vspace{9mm}
{\large 9 ноября 2026 г.\par}
\vspace{13mm}
{\normalsize Лёвин Артём Александрович\par}
{\normalsize эксклюзивно для Ксении Клыковой\par}
\vspace{13mm}
\begin{tikzpicture}[x=1cm,y=1cm]
 \draw[accent,line width=.95pt] (-4.0,0) .. controls (-2.4,.68) and (-1.3,-.56) .. (0,0)
 .. controls (1.5,.58) and (2.8,-.52) .. (4.0,0);
 \fill[accent] (0,0) circle (1.25pt);
\end{tikzpicture}
\vspace{9mm}
\end{titlepage}

\mainsection{1. Как устроен кредит с погашением по графику}
В задачах ЕГЭ часто задана \textbf{сумма кредита}, расписание \textbf{остатков долга} и правило начисления процентов. Нужно определить размер отдельных платежей, сумму выплат или допустимый размер кредита.

\subhead{Пять величин одной строки}
Для каждого периода удобно использовать одну и ту же последовательность:
\begin{center}
\begin{tikzpicture}[node distance=5mm, every node/.style={font=\small,align=center}, >=Latex]
\node (a) [draw=accent,rounded corners=2pt,inner sep=6pt,text width=25mm] {Начальный долг\\$B_{k-1}$};
\node (b) [draw=accent,rounded corners=2pt,inner sep=6pt,text width=25mm,right=of a] {Проценты\\$I_k=rB_{k-1}$};
\node (c) [draw=accent,rounded corners=2pt,inner sep=6pt,text width=27mm,right=of b] {Долг до выплаты\\$A_k=B_{k-1}+I_k$};
\node (d) [draw=accent,rounded corners=2pt,inner sep=6pt,text width=25mm,right=of c] {Платёж\\$P_k=A_k-B_k$};
\draw[->,darkaccent] (a)--(b);\draw[->,darkaccent](b)--(c);\draw[->,darkaccent](c)--(d);
\end{tikzpicture}
\end{center}
Здесь $r$~--- процентная ставка за \emph{один период начисления}, записанная десятичной дробью; $B_k$~--- долг после выплаты. Например, 12\% соответствует $r=0{,}12$.
\mathnote{\boxed{P_k=B_{k-1}(1+r)-B_k}}
\keyidea{Проценты начисляются на оставшуюся сумму долга. Последующие выплаты вычитаются после начисления процентов, если именно такой порядок указан в условии.}

\subhead{Что проверять в условии}
\begin{itemize}
\item \cbox Когда увеличивается долг: в январе, в конце месяца или в другой момент?
\item \cbox На какую сумму начисляют проценты: на первоначальный кредит или на остаток долга?
\item \cbox Когда вносят платежи и к какой дате относится таблица остатков?
\item \cbox В каких единицах заданы величины: рубли, тысячи или миллионы рублей?
\item \cbox Что именно ограничено: \emph{каждый} платёж или \emph{общая} сумма выплат?
\end{itemize}

\subhead{Небольшой пример}
Банк выдал 400 тыс. руб. На следующий январь долг увеличился на 10\%, а к июлю после выплат осталось 280 тыс. руб. Найдём, сколько заплатили за этот период:
\[
I=400\cdot0{,}10=40,\qquad A=400+40=440,\qquad
P=440-280=\boxed{160\text{ тыс. руб.}}
\]

\clearpage
\mainsection{2. Подробный разбор: кредит на 750 тысяч рублей}
В июле 2031 года взят кредит 750 тыс. руб. Каждый январь долг увеличивается на 12\%, а платежи вносятся с февраля по июнь. В июле 2032 и 2033 годов долг должен быть на 20\% меньше, чем в июле предыдущего года. К июлю 2034 года кредит нужно погасить полностью.

\textbf{Шаг 1.} Записываем остатки долга по июльским датам:
\[
B_0=750,\quad B_1=750\cdot0{,}8=600,\quad
B_2=600\cdot0{,}8=480,\quad B_3=0.
\]
\textbf{Шаг 2.} Для каждого года находим проценты, долг до погашения и величину платежа.

\begin{center}
\renewcommand{\arraystretch}{1.3}
\small
\begin{tabularx}{\linewidth}{@{}c r r r r@{}}
\toprule
Период & Начало, тыс. & Январские 12\% & Долг до выплат & Платёж, тыс.\\
\midrule
2032 & $750$ & $90$ & $840$ & $840-600=240$\\
2033 & $600$ & $72$ & $672$ & $672-480=192$\\
2034 & $480$ & $57{,}6$ & $537{,}6$ & $537{,}6-0=537{,}6$\\
\bottomrule
\end{tabularx}
\end{center}
\textbf{Шаг 3.} Складываем платежи:
\[
P_{\text{общ}}=240+192+537{,}6=\boxed{969{,}6\text{ тыс. руб.}}
\]
\textbf{Проверка:} общая переплата равна сумме всех процентов:
\[
90+72+57{,}6=219{,}6,\qquad 750+219{,}6=969{,}6.
\]
Таким способом можно быстро проверить большую таблицу без повторного сложения каждого платежа.

\subhead{Переплата в процентах}
Если $S$~--- сумма полученного кредита, а $P_{\text{общ}}$~--- все выплаты, то
\[
\text{переплата}=P_{\text{общ}}-S,\qquad
q=\frac{P_{\text{общ}}-S}{S}\cdot100\%.
\]
В разобранном примере $q=\dfrac{219{,}6}{750}\cdot100\%=\boxed{29{,}28\%}$.

\subhead{Типичная ошибка}
Нельзя вычислять 12\% от первоначальных 750 тыс. руб. ежегодно. Во втором году проценты берутся от 600 тыс., а в третьем~--- от 480 тыс. руб.

\clearpage
\mainsection{3. Максимальный кредит при ограничении каждого платежа}
Пусть $S$~--- положительное целое число миллионов рублей. В июле 2026 года взят кредит $S$. Каждый январь остаток увеличивается на 20\%, выплаты производятся с февраля по июнь, а долги в июле 2027, 2028 и 2029 годов должны составить $0{,}8S$, $0{,}4S$ и $0$ соответственно.

Найдём наибольшее целое $S$, при котором \textbf{каждый годовой платёж строго меньше 5 млн руб.}

\begin{center}
\small
\renewcommand{\arraystretch}{1.3}
\begin{tabularx}{\linewidth}{@{}c c c c c@{}}
\toprule
Год & До начисления & Проценты & До выплаты & Платёж\\
\midrule
2027 & $S$ & $0{,}2S$ & $1{,}2S$ & $0{,}4S$\\
2028 & $0{,}8S$ & $0{,}16S$ & $0{,}96S$ & $0{,}56S$\\
2029 & $0{,}4S$ & $0{,}08S$ & $0{,}48S$ & $0{,}48S$\\
\bottomrule
\end{tabularx}
\end{center}
Нужно одновременно выполнить три ограничения:
\[
\begin{cases}
0{,}4S<5,\\
0{,}56S<5,\\
0{,}48S<5.
\end{cases}
\quad\Longleftrightarrow\quad
\begin{cases}
S<12{,}5,\\[1mm]
S<\dfrac{125}{14}\approx8{,}93,\\[1mm]
S<\dfrac{125}{12}\approx10{,}42.
\end{cases}
\]
Наиболее строгое ограничение даёт \textbf{второй платёж}. Учитывая целочисленность $S$, получаем
\[
\boxed{S_{\max}=8\text{ млн руб.}}
\]
\keyidea{Фраза «каждый платёж меньше» означает систему неравенств. Достаточно найти наименьшую из верхних границ и проверить целочисленность и строгость знака.}

\subhead{Если ограничена вся сумма выплат}
По закону сохранения долга сумма всех выплат равна первоначальной сумме \textbf{плюс все начисленные проценты}:
\[
P_{\text{общ}}=S+\sum_{k=1}^{n}rB_{k-1}.
\]
Поэтому ограничение $P_{\text{общ}}<M$ приводит к \textbf{одному} неравенству. Его нельзя подменять ограничением каждого платежа.

\clearpage
\mainsection{4. Общая сумма выплат и проценты за несколько месяцев}

\subhead{Пример А. Четырёхлетний кредит под 15\%}
Размер кредита $S$ млн руб., где $S$~--- положительное целое число. Проценты начисляются ежегодно по ставке 15\%, а долги по июльским датам составляют $S$, $0{,}8S$, $0{,}5S$, $0{,}1S$, $0$. Общая сумма выплат должна быть \emph{меньше} 50 млн руб.

Находим проценты каждого года именно от остатка долга:
\[
I_1=0{,}15S,\quad I_2=0{,}12S,\quad
I_3=0{,}075S,\quad I_4=0{,}015S.
\]
Отсюда
\[
P_{\text{общ}}=S+(0{,}15+0{,}12+0{,}075+0{,}015)S
=1{,}36S.
\]
Следовательно,
\[
1{,}36S<50 \quad\Longleftrightarrow\quad
S<\frac{1250}{34}\approx36{,}76,
\qquad\boxed{S_{\max}=36}.
\]
\textbf{Проверка:} при $S=36$ общая сумма составит $48{,}96$ млн руб., а при $S=37$~--- $50{,}32$ млн руб., что уже превышает порог.

\subhead{Пример Б. Полугодовой кредит под 5\% в месяц}
15 января выдан кредит на сумму $S$. Долг на 15-е число каждого месяца, начиная с января и заканчивая июлем, задаётся в процентах от кредита: $100\%$, $90\%$, $80\%$, $70\%$, $60\%$, $50\%$, $0\%$. В конце каждого месяца (с января по июнь) долг увеличивается на 5\%, затем в первой половине следующего месяца вносится платёж.

На каждый из шести месяцев проценты начисляются от долга, который был \emph{до соответствующего начисления}:
\[
I_{\text{общ}}=0{,}05S(1+0{,}9+0{,}8+0{,}7+0{,}6+0{,}5)=0{,}225S.
\]
Общая сумма выплат $P_{\text{общ}}=1{,}225S$, поэтому переплата равна
\[
\frac{1{,}225S-S}{S}\cdot100\%=\boxed{22{,}5\%}.
\]
\keyidea{При ежемесячных процентах и меняющемся остатке нужно отдельно учесть проценты за каждый месяц.}

\clearpage
\thispagestyle{fancy}
\begin{center}
{\Large\bfseries\color{darkaccent}Алгоритм: кредит по графику остатков}\par
\vspace{5mm}
{\small Для задач с ежегодным или ежемесячным начислением процентов}\par
\end{center}
\vspace{7mm}
\begin{center}
\begin{tikzpicture}[>=Latex,node distance=10mm,
 act/.style={rectangle,rounded corners=3pt,draw=accent,line width=.8pt,text width=106mm,minimum height=13mm,align=center,inner sep=4pt,font=\small},
 decision/.style={diamond,aspect=3.7,draw=darkaccent,line width=.8pt,text width=65mm,align=center,inner sep=2pt,font=\small},
 ending/.style={ellipse,draw=accent,line width=.8pt,text width=70mm,minimum height=12mm,align=center,font=\small},
 ar/.style={->,darkaccent,line width=.95pt}]
\node (start) [ending] {Прочитайте условие и даты};
\node (dates) [act,below=of start] {Запишите остатки долга\\$B_0,B_1,\ldots,B_n$ и ставку $r$ за период};
\node (interest) [act,below=of dates] {Вычислите проценты каждого периода\\$I_k=rB_{k-1}$};
\node (payments) [act,below=of interest] {Найдите платежи\\$P_k=B_{k-1}+I_k-B_k$};
\node (what) [decision,below=11mm of payments] {Что требуется найти?};
\node (sum) [act,text width=51mm,anchor=north] at ([xshift=-40mm,yshift=-20mm]what.south) {Общая сумма:\\$S+\sum I_k$};
\node (bound) [act,text width=51mm,anchor=north] at ([xshift=40mm,yshift=-20mm]what.south) {Максимум $S$:\\составьте нужные неравенства};
\node (check) [act,text width=110mm,anchor=north] at ([yshift=-59mm]what.south) {Проверьте строгие знаки, единицы измерения,\\целочисленность и итоговый ответ};
\node (end) [ending,below=of check] {Ответ};
\draw[ar] (start)--(dates);\draw[ar](dates)--(interest);\draw[ar](interest)--(payments);\draw[ar](payments)--(what);
\draw[ar] (what.south west) -- node[above left,font=\scriptsize]{сумма} (sum.north);
\draw[ar] (what.south east) -- node[above right,font=\scriptsize]{граница} (bound.north);
\draw[ar] (sum.south) -- (check.north west);
\draw[ar] (bound.south) -- (check.north east);
\draw[ar] (check)--(end);
\end{tikzpicture}
\end{center}

\clearpage
\mainsection{5. Типичные ошибки и итоговая памятка}
\subhead{Четыре распространённые ошибки}
\begin{enumerate}[label=\arabic*.]
\item \textbf{Считают проценты от первоначального кредита каждый год.} При изменяющемся долге процент начисляется на \emph{остаток} соответствующего периода.
\item \textbf{Путают начисление и платёж.} Проценты сначала добавляются к долгу, затем из полученной суммы вычитается платёж.
\item \textbf{Меняют смысл ограничения.} Слова «каждый платёж» требуют отдельной проверки всех периодов; «общая сумма выплат»~--- сложения платежей или вычисления суммы начисленных процентов.
\item \textbf{Теряют строгость неравенства.} Если $P<5$, значение $P=5$ запрещено. Если сумма $S$ целая, нужно взять наибольшее \emph{допустимое целое} число.
\end{enumerate}

\subhead{Перед тем как записать ответ}
\begin{itemize}
\item \cbox Указаны все остатки долга в правильном порядке.
\item \cbox Начисленные проценты вычислены от нужного остатка.
\item \cbox Каждый платёж посчитан как разность долга до и после погашения.
\item \cbox Проверена формула $P_{\text{общ}}=S+I_{\text{общ}}$.
\item \cbox Все платежи неотрицательны, а последний остаток равен нулю.
\item \cbox Сохранены единицы измерения и знак «меньше» или «не больше».
\item \cbox При поиске максимального $S$ проверено соседнее большее целое число.
\end{itemize}

\subhead{Три формулы, которые стоит знать наизусть}
\[
\boxed{I_k=rB_{k-1}},\qquad
\boxed{P_k=B_{k-1}(1+r)-B_k},\qquad
\boxed{P_{\text{общ}}=S+\sum I_k}.
\]

\clearpage
\mainsection{6. Тренировочный блок}
Решите \textbf{10 задач} от базового уровня до повышенного. Все денежные величины считайте в указанных единицах. Если в условии сказано, что проценты начисляют в январе, а выплаты вносят с февраля по июнь, то долг на июль указан \emph{после} всех этих выплат. Для каждого задания записывайте исходный долг, начисленные проценты, долг до выплаты и итоговый долг.

\subhead{Базовый уровень}
\ex{1}{Один год начисления процентов}
В июле банк выдал клиенту кредит в размере 480 тысяч рублей. В январе следующего года долг по условиям договора увеличился на 10\% от оставшейся суммы. С февраля по июнь клиент погасил часть задолженности, и в июле долг составил ровно 300 тысяч рублей. Определите, сколько тысяч рублей клиент выплатил банку за период с февраля по июнь.

\ex{2}{Восстановление остатка долга}
Предприниматель получил в июле кредит на сумму 640 тысяч рублей. В январе следующего года банк единовременно начислил 12{,}5\% на всю оставшуюся задолженность. В период с февраля по июнь предприниматель внёс платежи на общую сумму 280 тысяч рублей. Найдите, какая сумма долга останется к июлю, после внесения всех указанных платежей. Ответ дайте в тысячах рублей.

\ex{3}{Два года и полное погашение}
В июле 2026 года был получен кредит в размере 500 тысяч рублей на два года. Каждый январь долг увеличивается на 10\% по сравнению с долгом на конец предыдущего года, а с февраля по июнь заёмщик выплачивает часть задолженности. Известно, что в июле 2027 года долг составил 300 тысяч рублей, а к июлю 2028 года кредит был полностью погашен. Найдите размер платежа за каждый год и общую сумму выплат. Все ответы запишите в тысячах рублей.

\subhead{Средний уровень}
\ex{4}{Три года: платежи и переплата}
В июле 2026 года оформлен кредит на сумму 900 тысяч рублей сроком на три года. В январе каждого года задолженность увеличивается на 10\%, а платежи осуществляются с февраля по июнь. По договору остаток долга в июле 2027 года должен составить 720 тысяч рублей, в июле 2028 года~--- 540 тысяч рублей, а в июле 2029 года~--- 0 рублей. Найдите все три годовых платежа, их общую сумму и процент переплаты относительно первоначально полученного кредита.

\ex{5}{График в долях от первоначального кредита}
Банк выдал кредит в размере 800 тысяч рублей на три года. Ежегодно в январе задолженность увеличивается на 15\%, а с февраля по июнь вносятся платежи. Остатки долга в июле первого, второго и третьего годов должны составлять соответственно 75\%, 50\% и 0\% от исходной суммы кредита. Определите размеры всех трёх платежей в тысячах рублей, общую сумму выплат и переплату в процентах от суммы кредита.

\vspace{6pt}
\subhead{Ограничения и поиск максимального кредита}
\ex{6}{Ограничение на каждый из двух платежей}
В июле предприниматель планирует взять в банке кредит на два года в размере $S$ миллионов рублей, где $S$~--- положительное целое число. Каждый январь долг увеличивается на 20\%; выплаты производятся с февраля по июнь. В июле первого года долг должен составлять 60\% от исходной суммы кредита, а в июле второго года кредит должен быть погашен. Найдите наибольшее целое значение $S$, если \textbf{каждый} из двух годовых платежей должен быть \textbf{строго меньше} 6 миллионов рублей.

\ex{7}{Какой платёж задаёт самое строгое ограничение?}
В июле 2026 года планируется кредит на три года в размере $S$ миллионов рублей, где $S$~--- положительное целое число. В январе каждого года банк начисляет 10\% на остаток долга; с февраля по июнь заёмщик выплачивает часть задолженности. По графику в июле 2027, 2028 и 2029 годов остатки должны составить $0{,}75S$, $0{,}35S$ и $0$ соответственно. Определите наибольшее целое $S$, если каждый отдельный годовой платёж должен быть строго меньше 4 миллионов рублей. Укажите, платёж за какой год создаёт самое жёсткое ограничение.

\ex{8}{Ограничение общей суммы выплат}
Банк готов выдать в июле кредит на четыре года в размере $S$ миллионов рублей, где $S$~--- положительное целое число. В январе каждого года оставшийся долг увеличивается на 12\%, а платежи вносятся с февраля по июнь. Долг по июльским датам, начиная с года выдачи кредита, должен последовательно составлять $S$, $0{,}85S$, $0{,}60S$, $0{,}25S$ и $0$. Найдите наибольшее целое значение $S$, при котором \textbf{общая сумма всех выплат} по кредиту будет строго меньше 30 миллионов рублей.

\vspace{6pt}
\subhead{Повышенный уровень}
\ex{9}{Ежемесячные начисления и процент переплаты}
15 января банк предоставил компании кредит на полгода в размере $S$ рублей. В конце каждого месяца, начиная с января и заканчивая июнем, задолженность увеличивается на 4\% от текущего остатка. В первой половине каждого месяца, начиная с февраля, компания вносит платёж; поэтому остатки долга на 15-е число месяцев с января по июль равны соответственно $100\%$, $90\%$, $80\%$, $65\%$, $45\%$, $20\%$ и $0\%$ от первоначального кредита. Определите, на сколько процентов общая сумма внесённых платежей превышает сумму полученного кредита. Покажите в решении, на какие остатки начислялись проценты.

\ex{10}{Два ограничения одновременно}
В июле 2026 года предприятие планирует взять кредит на три года в размере $S$ миллионов рублей, где $S$~--- положительное целое число. В январе каждого года задолженность увеличивается на 15\%, после чего с февраля по июнь предприятие вносит платежи. По графику остатки долга в июле 2027, 2028 и 2029 годов составят $0{,}7S$, $0{,}3S$ и $0$ соответственно. Банк требует, чтобы \textbf{каждый} годовой платёж был строго меньше 5 миллионов рублей, а \textbf{общая сумма всех выплат} была строго меньше 13 миллионов рублей. Найдите наибольшее целое значение $S$, удовлетворяющее обоим требованиям.

\clearpage
\mainsection{7. Ответы к тренировочному блоку}
Сначала выполните задания самостоятельно. Для задач с несколькими запрошенными величинами ответы приведены в том же порядке, что и в условии.

\vspace{5mm}
\renewcommand{\arraystretch}{1.22}
\begin{center}
\begin{tabularx}{\linewidth}{@{}c X@{}}
\toprule
\textbf{№} & \textbf{Правильный ответ}\\
\midrule
\ansrow{1}{228 тыс. руб.}
\ansrow{2}{440 тыс. руб.}
\ansrow{3}{250 тыс. руб.; 330 тыс. руб.; всего 580 тыс. руб.}
\ansrow{4}{270 тыс. руб.; 252 тыс. руб.; 594 тыс. руб.; всего 1116 тыс. руб.; переплата 24\%.}
\ansrow{5}{320 тыс. руб.; 290 тыс. руб.; 460 тыс. руб.; всего 1070 тыс. руб.; переплата 33{,}75\%.}
\ansrow{6}{$S_{\max}=8$ млн руб.}
\ansrow{7}{$S_{\max}=8$ млн руб.; самое строгое ограничение создаёт платёж за 2028 год.}
\ansrow{8}{$S_{\max}=22$ млн руб.}
\ansrow{9}{16\%.}
\ansrow{10}{$S_{\max}=9$ млн руб.}
\bottomrule
\end{tabularx}
\end{center}

\vspace{8mm}
\textcolor{darkaccent}{\bfseries Краткая самопроверка.}
Для задачи с платежами $P_1,\ldots,P_n$ и начальным кредитом $S$ обязательно должно выполняться
\[
\sum_{k=1}^nP_k=S+\sum_{k=1}^n I_k.
\]
Если равенство нарушено, проверьте порядок начисления процентов и вычитания платежей.

\vfill
\begin{center}
\color{softgray}\small Лёвин Артём Александрович\\
эксклюзивно для Ксении Клыковой
\end{center}
\end{document}
